%!TeX root=../gardentop.tex
\chapter{The Strangest House Any One Ever Lived In} 
\begin{figure}[t!]
\centering
\includegraphics[width=\linewidth]{august}
\end{figure}
 \lettrine[lines=6]{I}{t} was the sweetest, most mysterious-looking place any one could imagine. The high walls which shut it in were covered with the leafless stems of climbing roses which were so thick that they were matted together. Mary Lennox knew they were roses because she had seen a great many roses in India. All the ground was covered with grass of a wintry brown and out of it grew clumps of bushes which were surely rose-bushes if they were alive. There were numbers of standard roses which had so spread their branches that they were like little trees. There were other trees in the garden, and one of the things which made the place look strangest and loveliest was that climbing roses had run all over them and swung down long tendrils which made light swaying curtains, and here and there they had caught at each other or at a far-reaching branch and had crept from one tree to another and made lovely bridges of themselves. There were neither leaves nor roses on them now and Mary did not know whether they were dead or alive, but their thin gray or brown branches and sprays looked like a sort of hazy mantle spreading over everything, walls, and trees, and even brown grass, where they had fallen from their fastenings and run along the ground. It was this hazy tangle from tree to tree which made it all look so mysterious. Mary had thought it must be different from other gardens which had not been left all by themselves so long; and indeed it was different from any other place she had ever seen in her life.

»How still it is!« she whispered. »How still!«

Then she waited a moment and listened at the stillness. The robin, who had flown to his tree-top, was still as all the rest. He did not even flutter his wings; he sat without stirring, and looked at Mary.

»No wonder it is still,« she whispered again. »I am the first person who has spoken in here for ten years.«

She moved away from the door, stepping as softly as if she were afraid of awakening some one. She was glad that there was grass under her feet and that her steps made no sounds. She walked under one of the fairy-like gray arches between the trees and looked up at the sprays and tendrils which formed them.

»I wonder if they are all quite dead,« she said. »Is it all a quite dead garden? I wish it wasn't.«

If she had been Ben Weatherstaff she could have told whether the wood was alive by looking at it, but she could only see that there were only gray or brown sprays and branches and none showed any signs of even a tiny leaf-bud anywhere.

But she was inside the wonderful garden and she could come through the door under the ivy any time and she felt as if she had found a world all her own.

The sun was shining inside the four walls and the high arch of blue sky over this particular piece of Misselthwaite seemed even more brilliant and soft than it was over the moor. The robin flew down from his tree-top and hopped about or flew after her from one bush to another. He chirped a good deal and had a very busy air, as if he were showing her things. Everything was strange and silent and she seemed to be hundreds of miles away from any one, but somehow she did not feel lonely at all. All that troubled her was her wish that she knew whether all the roses were dead, or if perhaps some of them had lived and might put out leaves and buds as the weather got warmer. She did not want it to be a quite dead garden. If it were a quite alive garden, how wonderful it would be, and what thousands of roses would grow on every side!

Her skipping-rope had hung over her arm when she came in and after she had walked about for a while she thought she would skip round the whole garden, stopping when she wanted to look at things. There seemed to have been grass paths here and there, and in one or two corners there were alcoves of evergreen with stone seats or tall moss-covered flower urns in them.

As she came near the second of these alcoves she stopped skipping. There had once been a flower-bed in it, and she thought she saw something sticking out of the black earth—some sharp little pale green points. She remembered what Ben Weatherstaff had said and she knelt down to look at them.

»Yes, they are tiny growing things and they might be crocuses or snowdrops or daffodils,« she whispered.

She bent very close to them and sniffed the fresh scent of the damp earth. She liked it very much.

»Perhaps there are some other ones coming up in other places,« she said. »I will go all over the garden and look.«

She did not skip, but walked. She went slowly and kept her eyes on the ground. She looked in the old border beds and among the grass, and after she had gone round, trying to miss nothing, she had found ever so many more sharp, pale green points, and she had become quite excited again.

»It isn't a quite dead garden,« she cried out softly to herself. »Even if the roses are dead, there are other things alive.«

She did not know anything about gardening, but the grass seemed so thick in some of the places where the green points were pushing their way through that she thought they did not seem to have room enough to grow. She searched about until she found a rather sharp piece of wood and knelt down and dug and weeded out the weeds and grass until she made nice little clear places around them.

»Now they look as if they could breathe,« she said, after she had finished with the first ones. »I am going to do ever so many more. I'll do all I can see. If I haven't time to-day I can come to-morrow.«

She went from place to place, and dug and weeded, and enjoyed herself so immensely that she was led on from bed to bed and into the grass under the trees. The exercise made her so warm that she first threw her coat off, and then her hat, and without knowing it she was smiling down on to the grass and the pale green points all the time.

The robin was tremendously busy. He was very much pleased to see gardening begun on his own estate. He had often wondered at Ben Weatherstaff. Where gardening is done all sorts of delightful things to eat are turned up with the soil. Now here was this new kind of creature who was not half Ben's size and yet had had the sense to come into his garden and begin at once.

Mistress Mary worked in her garden until it was time to go to her midday dinner. In fact, she was rather late in remembering, and when she put on her coat and hat, and picked up her skipping-rope, she could not believe that she had been working two or three hours. She had been actually happy all the time; and dozens and dozens of the tiny, pale green points were to be seen in cleared places, looking twice as cheerful as they had looked before when the grass and weeds had been smothering them.

»I shall come back this afternoon,« she said, looking all round at her new kingdom, and speaking to the trees and the rose-bushes as if they heard her.

Then she ran lightly across the grass, pushed open the slow old door and slipped through it under the ivy. She had such red cheeks and such bright eyes and ate such a dinner that Martha was delighted.

»Two pieces o' meat an' two helps o' rice puddin'!« she said. »Eh! mother will be pleased when I tell her what th' skippin'-rope's done for thee.«

In the course of her digging with her pointed stick Mistress Mary had found herself digging up a sort of white root rather like an onion. She had put it back in its place and patted the earth carefully down on it and just now she wondered if Martha could tell her what it was.

»Martha,« she said, »what are those white roots that look like onions?«

»They're bulbs,« answered Martha. »Lots o' spring flowers grow from 'em. Th' very little ones are snowdrops an' crocuses an' th' big ones are narcissusis an' jonquils an' daffydowndillys. Th' biggest of all is lilies an' purple flags. Eh! they are nice. Dickon's got a whole lot of 'em planted in our bit o' garden.«

»Does Dickon know all about them?« asked Mary, a new idea taking possession of her.

»Our Dickon can make a flower grow out of a brick walk. Mother says he just whispers things out o' th' ground.«

»Do bulbs live a long time? Would they live years and years if no one helped them?« inquired Mary anxiously.

»They're things as helps themselves,« said Martha. »That's why poor folk can afford to have 'em. If you don't trouble 'em, most of 'em'll work away underground for a lifetime an' spread out an' have little 'uns. There's a place in th' park woods here where there's snowdrops by thousands. They're the prettiest sight in Yorkshire when th' spring comes. No one knows when they was first planted.«

»I wish the spring was here now,« said Mary. »I want to see all the things that grow in England.«

She had finished her dinner and gone to her favorite seat on the hearth-rug.

»I wish—I wish I had a little spade,« she said.

»Whatever does tha' want a spade for?« asked Martha, laughing. »Art tha' goin' to take to diggin'? I must tell mother that, too.«

Mary looked at the fire and pondered a little. She must be careful if she meant to keep her secret kingdom. She wasn't doing any harm, but if Mr~Craven found out about the open door he would be fearfully angry and get a new key and lock it up forevermore. She really could not bear that.

»This is such a big lonely place,« she said slowly, as if she were turning matters over in her mind. »The house is lonely, and the park is lonely, and the gardens are lonely. So many places seem shut up. I never did many things in India, but there were more people to look at—natives and soldiers marching by—and sometimes bands playing, and my Ayah told me stories. There is no one to talk to here except you and Ben Weatherstaff. And you have to do your work and Ben Weatherstaff won't speak to me often. I thought if I had a little spade I could dig somewhere as he does, and I might make a little garden if he would give me some seeds.«

Martha's face quite lighted up.

»There now!« she exclaimed, »if that wasn't one of th' things mother said. She says, 'There's such a lot o' room in that big place, why don't they give her a bit for herself, even if she doesn't plant nothin' but parsley an' radishes? She'd dig an' rake away an' be right down happy over it.' Them was the very words she said.«

»Were they?« said Mary. »How many things she knows, doesn't she?«

»Eh!« said Martha. »It's like she says: 'A woman as brings up twelve children learns something besides her A B C. Children's as good as 'rithmetic to set you findin' out things.'«

»How much would a spade cost—a little one?« Mary asked.

»Well,« was Martha's reflective answer, »at Thwaite village there's a shop or so an' I saw little garden sets with a spade an' a rake an' a fork all tied together for two shillings. An' they was stout enough to work with, too.«

»I've got more than that in my purse,« said Mary. »Mrs~Morrison gave me five shillings and Mrs~Medlock gave me some money from Mr~Craven.«

»Did he remember thee that much?« exclaimed Martha.

»Mrs~Medlock said I was to have a shilling a week to spend. She gives me one every Saturday. I didn't know what to spend it on.«

»My word! that's riches,« said Martha. »Tha' can buy anything in th' world tha' wants. Th' rent of our cottage is only one an' threepence an' it's like pullin' eye-teeth to get it. Now I've just thought of somethin',« putting her hands on her hips.

»What?« said Mary eagerly.

»In the shop at Thwaite they sell packages o' flower-seeds for a penny each, and our Dickon he knows which is th' prettiest ones an' how to make 'em grow. He walks over to Thwaite many a day just for th' fun of it. Does tha' know how to print letters?« suddenly.

»I know how to write,« Mary answered.

Martha shook her head.

»Our Dickon can only read printin'. If tha' could print we could write a letter to him an' ask him to go an' buy th' garden tools an' th' seeds at th' same time.«

»Oh! you're a good girl!« Mary cried. »You are, really! I didn't know you were so nice. I know I can print letters if I try. Let's ask Mrs~Medlock for a pen and ink and some paper.«

»I've got some of my own,« said Martha. »I bought 'em so I could print a bit of a letter to mother of a Sunday. I'll go and get it.«

She ran out of the room, and Mary stood by the fire and twisted her thin little hands together with sheer pleasure.

»If I have a spade,« she whispered, »I can make the earth nice and soft and dig up weeds. If I have seeds and can make flowers grow the garden won't be dead at all—it will come alive.«

She did not go out again that afternoon because when Martha returned with her pen and ink and paper she was obliged to clear the table and carry the plates and dishes down-stairs and when she got into the kitchen Mrs~Medlock was there and told her to do something, so Mary waited for what seemed to her a long time before she came back. Then it was a serious piece of work to write to Dickon. Mary had been taught very little because her governesses had disliked her too much to stay with her. She could not spell particularly well but she found that she could print letters when she tried. This was the letter Martha dictated to her:

\begin{quotation}
\noindent My Dear Dickon:

This comes hoping to find you well as it leaves me at present. Miss Mary has plenty of money and will you go to Thwaite and buy her some flower seeds and a set of garden tools to make a flower-bed. Pick the prettiest ones and easy to grow because she has never done it before and lived in India which is different. Give my love to mother and every one of you. Miss Mary is going to tell me a lot more so that on my next day out you can hear about elephants and camels and gentlemen going hunting lions and tigers.

\begin{flushright}
Your loving sister,\\
Martha Phœbe Sowerby.
\end{flushright}
\end{quotation}

»We'll put the money in th' envelope an' I'll get th' butcher's boy to take it in his cart. He's a great friend o' Dickon's,« said Martha.

»How shall I get the things when Dickon buys them?« asked Mary.

»He'll bring 'em to you himself. He'll like to walk over this way.«

»Oh!« exclaimed Mary, »then I shall see him! I never thought I should see Dickon.«

»Does tha' want to see him?« asked Martha suddenly, she had looked so pleased.

»Yes, I do. I never saw a boy foxes and crows loved. I want to see him very much.«

Martha gave a little start, as if she suddenly remembered something.

»Now to think,« she broke out, »to think o' me forgettin' that there; an' I thought I was goin' to tell you first thing this mornin'. I asked mother—and she said she'd ask Mrs~Medlock her own self.«

»Do you mean—« Mary began.

»What I said Tuesday. Ask her if you might be driven over to our cottage some day and have a bit o' mother's hot oat cake, an' butter, an' a glass o' milk.«

It seemed as if all the interesting things were happening in one day. To think of going over the moor in the daylight and when the sky was blue! To think of going into the cottage which held twelve children!

»Does she think Mrs~Medlock would let me go?« she asked, quite anxiously.

»Aye, she thinks she would. She knows what a tidy woman mother is and how clean she keeps the cottage.«

»If I went I should see your mother as well as Dickon,« said Mary, thinking it over and liking the idea very much. »She doesn't seem to be like the mothers in India.«

Her work in the garden and the excitement of the afternoon ended by making her feel quiet and thoughtful. Martha stayed with her until tea-time, but they sat in comfortable quiet and talked very little. But just before Martha went down-stairs for the tea-tray, Mary asked a question.

»Martha,« she said, »has the scullery-maid had the toothache again to-day?«

Martha certainly started slightly.

»What makes thee ask that?« she said.

»Because when I waited so long for you to come back I opened the door and walked down the corridor to see if you were coming. And I heard that far-off crying again, just as we heard it the other night. There isn't a wind to-day, so you see it couldn't have been the wind.«

»Eh!« said Martha restlessly. »Tha' mustn't go walkin' about in corridors an' listenin'. Mr~Craven would be that there angry there's no knowin' what he'd do.«

»I wasn't listening,« said Mary. »I was just waiting for you—and I heard it. That's three times.«

»My word! There's Mrs~Medlock's bell,« said Martha, and she almost ran out of the room.

»It's the strangest house any one ever lived in,« said Mary drowsily, as she dropped her head on the cushioned seat of the armchair near her. Fresh air, and digging, and skipping-rope had made her feel so comfortably tired that she fell asleep.